\section{Symmetry in Quantum Mechanics}

\subsection{Symmetries, Conservation Laws, and Degeneracies}
\subsubsection{Symmetries in Classical Physics}
We will review the concept of symmetry and conservation in classical physics. Suppose we have a Lagrangian $L$, a function of the generalized coordinate $q_i$ and the corresponding velocity $\dot{q}_i$. For a displacement
\begin{equation}
    q_i \to q_i + \delta q_i,
\end{equation}
if the Lagrangian remains unchanged, we have
\begin{equation}
    \frac{\partial L}{\partial q_i} = 0.
\end{equation}
From the Lagrange equation it then follows that 
\begin{equation}
    \frac{d p_i}{dt} = 0,
\end{equation}
where the canonical momentum is defined by
\begin{equation}
    p_i = \frac{\partial L}{\partial \dot{q}_i}.
\end{equation}
So the canonical momentum conjugate to $q_i$ is a conserved quantity under displacement if $L$ remains unchanged.

Similarly, looking at the Hamiltonian formulation we have $H$ as a function of $q_i$ and $p_i$, we have
\begin{equation}
    \frac{d p_i}{dt} = 0
\end{equation}
when 
\begin{equation}
    \frac{\partial H}{\partial t} = 0.
\end{equation}
Thus, if we lack an explicit $q_i$ dependence in $H$, ($H$ has a symmetry under $q_i \to q_i + \delta q_i$) we have a conserved quantity.

\subsubsection{Symmetry in Quantum Mechanics}
In quantum mechanics, we have made the association between an operation, like translation or rotation, to a unitary operator, like $\mathscr{S}$. Regardless if the physical system has a symmetry corresponding to the operator $\mathscr{S}$, we tend to call this operator for a symmetry operator. Furthermore, we have seen that if the symmetry operation is infinitesimal we can write
\begin{equation}
    \mathscr{S} = 1 - \frac{i \epsilon}{\hbar} G,
\end{equation}
where $G$ is the Hermitian generator of the symmetry operator. 

Consider that $H$ is invariant under $\mathscr{S}$. From this, it follows that
\begin{equation}
    \mathscr{S}^\dagger H \mathscr{S} = H,
\end{equation}
or equivalently
\begin{equation}
    \comm{G}{H} = 0.
\end{equation}
From Heisenberg's equation of motion ($i\hbar \partial_t A = \comm{A}{H}$),  we obtain
\begin{equation}
    \frac{d G}{dt} = 0,
\end{equation}
and thus, $G$ is a constant of the motion. E.g. if the Hamiltonian is translationally invariant, the momentum (that generates the translation), is a constant of the motion. 

Consider the system in an eigenket of $G$, when $G$ commutes with $H$, at time $t_0$. For a later time $t>t_0$, we have the ket
\begin{equation}
    \k{g',t_0;t} = U(t,t_0) \k{g'},
\end{equation}
which is also an eigenket of $G$ the same eigenvalue $g'$. Thus, when a ket is an eigenket to $G$ with eigenvalue $g'$ it will stay as such. Proving this is just realizing that since $\comm{G}{H} = 0$, then $\comm{G}{U} = 0$.

\subsubsection{Degeneracies}
Suppose that for a symmetry operator $\mathscr{S}$ we have
\begin{equation}
    \comm{H}{\mathscr{S}} = 0,
\end{equation}
and $\k{n}$ is an energy eigenket with eigenvalue $E_n$. Evidently, $\mathscr{S}\k{n}$ is also an energy eigenket with the same energy. If $\k{n}$ and $\mathscr{S}\k{n}$ represent different states we say that the states are degenerate.  Often, the symmetry operator is a function of a continuous parameter, say $\lambda$, and thus all states of the form $\mathscr{S}(\lambda)\k{n}$, are degenerate.

Consider that the symmetry operator in question is rotation, that is we have $\mathscr{D}(R)$. We also suppose that the Hamiltonian is rotationally invariant
\begin{equation}
    \comm{\mathscr{D}(R)}{H} = 0.
\end{equation}
This implies the commutation relations
\begin{alphaAlign}
    \comm{\vec{J}}{H} &= 0\\
    \comm{\vec{J}^2}{H} &= 0,
\end{alphaAlign}
and we can have simultaneous eigenkets of $H$, $\vec{J}^2$, and $J_z$, denoted by $\k{n; j, m}$. All states on the form
\begin{equation}
    \mathscr{R} \k{n; j, m}
\end{equation}
are degenerate and share the same energy. During rotation, different $m$ values are mixed, such that generally, $\mathscr{D}(R)\k{n; j, m}$ is a linear combination of $2j +1$ independent states. We can write this as
\begin{equation}
    \mathscr{D}(R)\k{n; j, m} = \sum_{m'} \k{n; j, m'} \mathscr{D}^{(j)}_{m'm}(R).
\end{equation}
This degeneracy is equal to the number of $m$-values, that is $2j+1$. It is also possible to reach this conclusion by applying $J_\pm$ to the ket.

As an example of this consider an atomic electron subject to the potential $V = V(r) + V_{LS}(r) \vec{L} \cdot \vec{S}$. Since this potential is rotationally invariant, we expect a $(2j+1)$-fold degeneracy. If we applied an external magnetic field, this symmetry would break, and we would get distinct energy levels. 

\subsubsection{SO(4) Symmetry in the Coulomb Potential}
An archetypal example of continuous symmetry in quantum mechanics is the solution of the coulomb potential in the hydrogen atom. The energy eigenvalue for this potential is
\begin{equation}
    E_n = \frac{1}{2} mc^2 \frac{Z^2 \alpha^2}{n^2} = \qty{-13.6}{\electronvolt} \frac{Z^2}{n^2},
\end{equation}
and since this energy is only dependent on $n$, the degeneracy for the ket $\k{n,l,m}$ is 
\begin{equation}
    \sum_{l=0}^{n-1}(2l+1) = n^2.
\end{equation}
This is due to the symmetry of $1/r$ potentials.

Classically, the orbits of $1/r$ potentials have been studied, namely the Kepler problem. The result from this leads us to the conclusion of the existence of a constant of the motion, which in classical mechanics is
\begin{equation}
    \vec{M} = \frac{\vec{p}\times\vec{L}}{m} - \frac{Ze^2}{r}\vec{x},
\end{equation}
called the Lenz vector, or the Runge-Lenz vector.

Treating this quantum mechanically, this new symmetry is called SO(4), analogous to the SO(3) group, but instead rotation operators in four spatial dimensions. We start by modifying the Lenz vector such that it becomes Hermitian
\begin{equation}
    \vec{M} = \frac{1}{2m} (\vec{p}\times\vec{L} - \vec{L}\times\vec{p}) - \frac{Z e^2}{r}\vec{x}.
\end{equation}
We then have
\begin{equation}
    \comm{\vec{M}}{H} = 0,
\end{equation}
for the Hamiltonian
\begin{equation}
    H = \frac{\vec{p}^2}{2m} - \frac{Ze^2}{r},
\end{equation}
and thus $\vec{M}$ is a constant of the motion. Furthermore, the relations
\begin{align}
    \vec{L}\cdot\vec{M} &= \vec{M}\cdot\vec{L} = 0,\\
    \vec{M}^2 &= \frac{2}{m} H (\vec{L}^2 + \hbar^2) + Z^2 e^4, \label{eq:M2}
\end{align}
are useful.

We want to identify the responsible symmetry for this constant of the motion. This is done by exploring the algebra of the generators of this symmetry. We know
\begin{equation}
    \comm{L_i}{L_j} = i\hbar \epsilon_{ijk} L_k,
\end{equation}
and it is possible to show
\begin{equation}
    \comm{M_i}{L_j} = i\hbar \epsilon_{ijk} M_k.
\end{equation}
With this we can derive that
\begin{equation}
    \comm{M_i}{M_j} = -i\hbar \epsilon_{ijk} \frac{2}{m} H L_k.
\end{equation}
These commutation relations do not form a closed algebra as a consequence of the presence of $H$, making it challenging identifying these operators as the generators. Conversely, we will consider the problem with specific bound states. The vector space will now only contain the eigenstates to $H$ that have an energy eigenvalue of $E < 0$, and we may replace $H$ with $E$, closing the algebra. We introduce a scaled vector operator 
\begin{equation}
    \vec{N} = \sqrt{-\frac{m}{2E}} \vec{M},
\end{equation}
such that the closed algebra reads as
\begin{alphaAlign}
    \comm{L_i}{L_j} &= i\hbar \epsilon_{ijk} L_k\\
    \comm{N_i}{L_j} &= i\hbar \epsilon_{ijk} N_k\\
    \comm{N_i}{N_j} &= i\hbar \epsilon_{ijk} L_k,
\end{alphaAlign}

The question becomes what type of symmetry is generated by the operators $\vec{L}$ and $\vec{N}$? Interestingly, the answer is rotation in four spatial dimensions. The first clue for this is that the required number of operators for rotation in $n$ dimensions is given by $n(n-1)/2$, and thus rotation in four dimensions require six operators, which we have with $\vec{L}$ and $\vec{n}$. It is less obvious to see that the above algebra is appropriate for this.

We generalize the procedure performed for three dimensions with $\vec{L}$.  We start by rewriting $(x,y,z)$ and $(p_x, p_y, p_z)$ to $(x_1,x_2, x_3)$ and $(p_1, p_2, p_3)$. Consequently, we write
\begin{alphaAlign}
    L_1 = \tilde{L}_{23} &= x_2 p_3 - x_3 p_2,\\
    L_2 = \tilde{L}_{31} &= x_3 p_1 - x_1 p_3,\\
    L_3 = \tilde{L}_{12} &= x_1 p_2 - x_2 p_1.
\end{alphaAlign}
We know invent a fourth spatial dimension $x_4$ and define its conjugate momentum $p_4$ according to the canonical commutation relation, and write
\begin{alphaAlign}
    N_1 = \tilde{L}_{14} &= x_1 p_4 - x_4 p_1,\\
    N_2 = \tilde{L}_{24} &= x_2 p_4 - x_4 p_2,\\
    N_3 = \tilde{L}_{34} &= x_3 p_4 - x_4 p_3.
\end{alphaAlign}
These operators obey the closed algebra defined earlier. That is, the closed algebra we found, is the algebra for rotation in four spatial dimensions.

We continue by investigating the degeneracies in the Coulomb potential. Starting by defining the operators
\begin{align}
    \vec{I} &= \frac{\vec{L} + \vec{N}}{2},\\
    \vec{K} &= \frac{\vec{L} - \vec{N}}{2},
\end{align}
we can prove the algebra
\begin{alphaAlign}
    \comm{I_i}{I_j} &= i\hbar \epsilon_{ijk} I_k\\
    \comm{K_i}{K_j} &= i\hbar \epsilon_{ijk} K_k\\
    \comm{I_i}{K_j} &= 0.
\end{alphaAlign}
As such, these operators obey independent angular momentum algebra, and we also have that $\comm{\vec{I}}{I} = \comm{\vec{K}}{H} = 0$. We denote the eigenvalues of $\vec{I}^2$ and $\vec{K}^2$ by $i(i+1)\hbar^2$ and $k(k+1)\hbar^2$ respectively, with $i,k$ taking on positive half-integer values.

We find that 
\begin{alphaAlign}
    \vec{I}^2 - \vec{K}^2 &= \frac{1}{4} \ab[ (\vec{L} + \vec{N})^2 - (\vec{L} - \vec{N})^2 ]\\
    &= \frac{1}{4} \ab[ (\vec{L}^2 + \vec{N}^2 + 2\vec{L}\cdot\vec{N}) - (\vec{L}^2 + \vec{N}^2 - 2\vec{L}\cdot\vec{N}) ]\\
    &= \vec{L} \cdot \vec{N}\\
    &= 0,
\end{alphaAlign}
and thus we must have that $i=k$. Similarly, we find
\begin{alphaAlign}
    \vec{I}^2 + \vec{K}^2 &= \frac{1}{4} \ab[ (\vec{L} + \vec{N})^2 + (\vec{L} - \vec{N})^2 ]\\
    &= \frac{1}{4} \ab[ (\vec{L}^2 + \vec{N}^2 + 2\vec{L}\cdot\vec{N}) + (\vec{L}^2 + \vec{N}^2 - 2\vec{L}\cdot\vec{N}) ]\\
    &= \frac{1}{2} \ab [ \vec{L}^2 + \vec{N}^2 ]\\
    &= \frac{1}{2} \ab [ \vec{L}^2 - \frac{m}{2E} \vec{M}^2 ],
\end{alphaAlign}
that together with Eq. \eqref{eq:M2} gives
\begin{equation}
    2k(k+1)\hbar^2 = \frac{1}{2} \ab ( -\hbar^2 - \frac{m}{2E}  Z^2 e^4).
\end{equation}
Solving this for $E$ yields the relation
\begin{equation}
    E = -\frac{mZ^2 e^4}{2\hbar^2} \frac{1}{(2k+1)^2}.
\end{equation}
This is the energy for the hydrogen atom but with $n$ replaced with $2k+1$. Thus, the rotational symmetries that arise from $\vec{I}$ and $\vec{K}$ creates the degeneracy seen. Technically the degeneracy is $(2i+1)(2k+1) = (2k+1)^2 = n^2$. 

Note here that we found the eigenvalues of the hydrogen atom without solving the Schrödinger equation. Instead, we used the symmetries in the system to obtain the same answer.

The SO(4) algebra used here can also, as can be seen from the above calculations be thought of as SU(2)$\times$SU(2), as $\vec{I}$ and $\vec{K}$ independently follow the SU(2) group algebra.





